\begin{table}[!htb]
\centering
\footnotesize
\setlength{\tabcolsep}{5pt}
\begin{tabular}{llcccc}
\toprule
& & \multicolumn{2}{c}{$C_A$} & \multicolumn{2}{c}{$T_Fn_f$}\\
\cmidrule(lr){3-4}\cmidrule(lr){5-6}
& & $HH$ & $AA$ & $HH$ & $AA$\\
\midrule
\multicolumn{6}{l}{$\nu=0$}\\
$x=\frac58$ & amplitudes & $4.96102843$ & $4.99578015$ & $-4.72601066$ & $-4.79551411$\\
 & closed form & $4.96102843$ & $4.99578015$ & $-4.72601066$ & $-4.79551411$\\
$x=\frac34$ & amplitudes & $5.18575874$ & $5.20342420$ & $-4.67802109$ & $-4.71335201$\\
 & closed form & $5.18575874$ & $5.20342420$ & $-4.67802109$ & $-4.71335201$\\
$x=\frac56$ & amplitudes & $5.44666610$ & $5.45382124$ & $-4.62275054$ & $-4.63706082$\\
 & closed form & $5.44666610$ & $5.45382124$ & $-4.62275054$ & $-4.63706082$\\
\multicolumn{2}{l}{largest difference} & $5.4\times10^{-10}$ & $6.4\times10^{-10}$ & $2.1\times10^{-10}$ & $2.5\times10^{-10}$\\
\multicolumn{2}{l}{window shift} & $4.1\times10^{-8}$ & $4.9\times10^{-8}$ & $1.7\times10^{-8}$ & $2.0\times10^{-8}$\\
\midrule
\multicolumn{6}{l}{$\nu=2$}\\
$x=\frac58$ & amplitudes & $5.72676090$ & $3.91372974$ & $-6.25747561$ & $-2.63141328$\\
 & closed form & $5.72676090$ & $3.91372974$ & $-6.25747561$ & $-2.63141328$\\
$x=\frac34$ & amplitudes & $5.45921432$ & $4.67872651$ & $-5.22493225$ & $-3.66395664$\\
 & closed form & $5.45921432$ & $4.67872651$ & $-5.22493225$ & $-3.66395664$\\
$x=\frac56$ & amplitudes & $5.54920634$ & $5.16581976$ & $-4.82783103$ & $-4.06105786$\\
 & closed form & $5.54920634$ & $5.16581976$ & $-4.82783103$ & $-4.06105786$\\
\multicolumn{2}{l}{largest difference} & $1.5\times10^{-10}$ & $1.0\times10^{-10}$ & $5.6\times10^{-11}$ & $6.8\times10^{-11}$\\
\multicolumn{2}{l}{window shift} & $6.6\times10^{-9}$ & $4.4\times10^{-9}$ & $4.5\times10^{-9}$ & $5.5\times10^{-9}$\\
\bottomrule
\end{tabular}
\caption{The constant $r_0$ of the one-loop kernel read off the two-loop $gg\to HH$ and $gg\to AA$ amplitudes with Eq.~(\ref{eq:r0read}), for the two components and the two colour factors of gluon splittings, at the three velocities of the scans, against the closed forms: Eq.~(\Lref{eq:r0g}) with $F_0=F_S$ for $HH$, and Eq.~(\ref{eq:r0aa}) for $AA$, which has $F_0=F_A$ for $\nu=0$ and the helicity-flip term with the opposite sign for $\nu=2$. The two processes differ by the helicity term alone in both components. The rows marked largest difference give, for each component, the largest difference between the amplitudes and the closed forms over the three velocities, and the rows marked window shift the largest change of the fitted value when the fit window is moved one step up in $\delta$ (Sec.~\ref{app:fits}).}
\label{tab:r0amp}
\end{table}
